%File: formatting-instructions-latex-2026.tex
%release 2026.0
\pdfoutput=1
\documentclass[letterpaper]{article} % DO NOT CHANGE THIS
\usepackage{aaai2026}  % DO NOT CHANGE THIS
\usepackage{times}  % DO NOT CHANGE THIS
\usepackage{helvet}  % DO NOT CHANGE THIS
\usepackage{courier}  % DO NOT CHANGE THIS
\usepackage[hyphens]{url}  % DO NOT CHANGE THIS
\usepackage{graphicx} % DO NOT CHANGE THIS
\urlstyle{rm} % DO NOT CHANGE THIS
\def\UrlFont{\rm}  % DO NOT CHANGE THIS
\usepackage{natbib}  % DO NOT CHANGE THIS AND DO NOT ADD ANY OPTIONS TO IT
\usepackage{caption} % DO NOT CHANGE THIS AND DO NOT ADD ANY OPTIONS TO IT
\frenchspacing  % DO NOT CHANGE THIS
\setlength{\pdfpagewidth}{8.5in}  % DO NOT CHANGE THIS
\setlength{\pdfpageheight}{11in}  % DO NOT CHANGE THIS
%
% These are recommended to typeset algorithms but not required. See the subsubsection on algorithms. Remove them if you don't have algorithms in your paper.
\usepackage{algorithm}
\usepackage{algorithmic}

%
% These are are recommended to typeset listings but not required. See the subsubsection on listing. Remove this block if you don't have listings in your paper.
\usepackage{newfloat}
\usepackage{listings}
\DeclareCaptionStyle{ruled}{labelfont=normalfont,labelsep=colon,strut=off} % DO NOT CHANGE THIS
\lstset{%
	basicstyle={\footnotesize\ttfamily},% footnotesize acceptable for monospace
	numbers=left,numberstyle=\footnotesize,xleftmargin=2em,% show line numbers, remove this entire line if you don't want the numbers.
	aboveskip=0pt,belowskip=0pt,%
	showstringspaces=false,tabsize=2,breaklines=true}
\floatstyle{ruled}
\newfloat{listing}{tb}{lst}{}
\floatname{listing}{Listing}
%
% Keep the \pdfinfo as shown here. There's no need
% for you to add the /Title and /Author tags.
\pdfinfo{
/TemplateVersion (2026.1)
}

\usepackage{amsmath}
\usepackage{pifont}
\usepackage{booktabs}
\usepackage{amsfonts}
\usepackage{hyperref}
\usepackage{cleveref}
\crefname{figure}{Fig.}{Figs.}
\crefname{table}{Table.}{Tables.}
\crefname{section}{Section}{Secs.}

% DISALLOWED PACKAGES
% \usepackage{authblk} -- This package is specifically forbidden
% \usepackage{balance} -- This package is specifically forbidden
% \usepackage{color (if used in text)
% \usepackage{CJK} -- This package is specifically forbidden
% \usepackage{float} -- This package is specifically forbidden
% \usepackage{flushend} -- This package is specifically forbidden
% \usepackage{fontenc} -- This package is specifically forbidden
% \usepackage{fullpage} -- This package is specifically forbidden
% \usepackage{geometry} -- This package is specifically forbidden
% \usepackage{grffile} -- This package is specifically forbidden
% \usepackage{hyperref} -- This package is specifically forbidden
% \usepackage{navigator} -- This package is specifically forbidden
% (or any other package that embeds links such as navigator or hyperref)
% \indentfirst} -- This package is specifically forbidden
% \layout} -- This package is specifically forbidden
% \multicol} -- This package is specifically forbidden
% \nameref} -- This package is specifically forbidden
% \usepackage{savetrees} -- This package is specifically forbidden
% \usepackage{setspace} -- This package is specifically forbidden
% \usepackage{stfloats} -- This package is specifically forbidden
% \usepackage{tabu} -- This package is specifically forbidden
% \usepackage{titlesec} -- This package is specifically forbidden
% \usepackage{tocbibind} -- This package is specifically forbidden
% \usepackage{ulem} -- This package is specifically forbidden
% \usepackage{wrapfig} -- This package is specifically forbidden
% DISALLOWED COMMANDS
% \nocopyright -- Your paper will not be published if you use this command
% \addtolength -- This command may not be used
% \balance -- This command may not be used
% \baselinestretch -- Your paper will not be published if you use this command
% \clearpage -- No page breaks of any kind may be used for the final version of your paper
% \columnsep -- This command may not be used
% \newpage -- No page breaks of any kind may be used for the final version of your paper
% \pagebreak -- No page breaks of any kind may be used for the final version of your paperr
% \pagestyle -- This command may not be used
% \tiny -- This is not an acceptable font size.
% \vspace{- -- No negative value may be used in proximity of a caption, figure, table, section, subsection, subsubsection, or reference
% \vskip{- -- No negative value may be used to alter spacing above or below a caption, figure, table, section, subsection, subsubsection, or reference

\setcounter{secnumdepth}{2} %May be changed to 1 or 2 if section numbers are desired.

% The file aaai2026.sty is the style file for AAAI Press
% proceedings, working notes, and technical reports.
%

% Title

% Your title must be in mixed case, not sentence case.
% That means all verbs (including short verbs like be, is, using,and go),
% nouns, adverbs, adjectives should be capitalized, including both words in hyphenated terms, while
% articles, conjunctions, and prepositions are lower case unless they
% directly follow a colon or long dash
\iffalse
\title{D-FCGS: Feedforward Compression of Dynamic Gaussian Splatting for Free-Viewpoint Videos}
\author{
    %Authors
    % All authors must be in the same font size and format.
    Written by AAAI Press Staff\textsuperscript{\rm 1}\thanks{With help from the AAAI Publications Committee.}\\
    AAAI Style Contributions by Pater Patel Schneider,
    Sunil Issar,\\
    J. Scott Penberthy,
    George Ferguson,
    Hans Guesgen,
    Francisco Cruz\equalcontrib,
    Marc Pujol-Gonzalez\equalcontrib
}
\affiliations{
    %Afiliations
    \textsuperscript{\rm 1}Association for the Advancement of Artificial Intelligence\\
    % If you have multiple authors and multiple affiliations
    % use superscripts in text and roman font to identify them.
    % For example,

    % Sunil Issar\textsuperscript{\rm 2}, 
    % J. Scott Penberthy\textsuperscript{\rm 3}, 
    % George Ferguson\textsuperscript{\rm 4},
    % Hans Guesgen\textsuperscript{\rm 5}
    % Note that the comma should be placed after the superscript

    1101 Pennsylvania Ave, NW Suite 300\\
    Washington, DC 20004 USA\\
    % email address must be in roman text type, not monospace or sans serif
    proceedings-questions@aaai.org
%
% See more examples next
}
\fi

%Example, Single Author, ->> remove \iffalse,\fi and place them surrounding AAAI title to use it
\iffalse
\title{My Publication Title --- Single Author}
\author {
    Author Name
}
\affiliations{
    Affiliation\\
    Affiliation Line 2\\
    name@example.com
}
\fi

%Example, Multiple Authors, ->> remove \iffalse,\fi and place them surrounding AAAI title to use it
\title{D-FCGS: Feedforward Compression of Dynamic Gaussian Splatting for Free-Viewpoint Videos}
\author {
    % Authors
    Wenkang Zhang\textsuperscript{\rm 1, \rm 2},
    Yan Zhao\textsuperscript{\rm 1},
    Qiang Wang\textsuperscript{\rm 3},
    Zhixin Xu\textsuperscript{\rm 1}, 
    Li Song\textsuperscript{\rm 1}, 
    Zhengxue Cheng\textsuperscript{\rm 1}\thanks{Corresponding Author}
}
\affiliations {
    % Affiliations
    \textsuperscript{\rm 1}Shanghai Jiao Tong University\\
    \textsuperscript{\rm 2}Zhiyuan College, Shanghai Jiao Tong University\\
    \textsuperscript{\rm 3}Visionstar Information Technology (Shanghai) Co., Ltd\\
    \{conquer.wkzhang, zhaoyanzy, zhixin.xu, song\_li, zxcheng\}@sjtu.edu.cn, wq@sightp.com
}


% REMOVE THIS: bibentry
% This is only needed to show inline citations in the guidelines document. You should not need it and can safely delete it.
\usepackage{bibentry}
% END REMOVE bibentry

\begin{document}

\maketitle

\begin{abstract}
Free-Viewpoint Video (FVV) enables immersive 3D experiences, but efficient compression of dynamic 3D representation remains a major challenge.
%
Existing dynamic 3D Gaussian Splatting methods couple reconstruction with optimization-dependent compression and customized motion formats, limiting generalization and standardization. 
%
To address this, we propose D-FCGS, a novel \textbf{F}eedforward \textbf{C}ompression framework for \textbf{D}ynamic \textbf{G}aussian \textbf{S}platting. Key innovations include: (1) a standardized Group-of-Frames (GoF) structure with I-P coding, leveraging sparse control points to extract inter-frame motion tensors; (2) a dual prior-aware entropy model that fuses hyperprior and spatial-temporal priors for accurate rate estimation; (3) a control-point-guided motion compensation mechanism and refinement network to enhance view-consistent fidelity. 
%
Trained on Gaussian frames derived from multi-view videos, D-FCGS generalizes across diverse scenes in a zero-shot fashion. Experiments show that it matches the rate-distortion performance of optimization-based methods, achieving over 17$\times$ compression compared to the baseline while preserving visual quality across viewpoints. This work advances feedforward compression of dynamic 3DGS, facilitating scalable FVV transmission and storage for immersive applications.
\end{abstract}

% Uncomment the following to link to your code, datasets, an extended version or similar.
% You must keep this block between (not within) the abstract and the main body of the paper.
\begin{links}
    \link{Code}{https://github.com/Mr-Zwkid/D-FCGS}
    % \link{Datasets}{https://aaai.org/example/datasets}
    % \link{Extended version}{https://arxiv.org/abs/2507.05859}
\end{links}

\section{Introduction}
Our world is inherently dynamic, with 3D scenes evolving over time and observable from arbitrary viewpoints. Capturing and representing such 4D dynamics has long been a fundamental challenge in computer vision and graphics. Free-viewpoint video (FVV), which enables immersive 6-DoF experiences, has emerged as a promising solution with applications in virtual reality, telepresence, and remote education. However, realizing practical FVV systems demands efficient solutions for reconstruction, compression, transmission, and rendering. In this work, we focus on the critical problem of compression for dynamic 3D scenes.

Recent advances of 3D Gaussian Splatting (3DGS) \cite{kerbl20233d} have revolutionized 3D scene representation, offering unparalleled rendering quality and real-time performance. Extending to 4D, dynamic forms of 3DGS \cite{yang2024deformable, wu20244d, li2024spacetime, yang2023real, sun20243dgstream} gradually garner attention. Analogous to the temporal expansion of images facilitated by videos, frame-by-frame 3D Gaussians can naturally serve as a temporal expansion of 3DGS, forming the basis for FVV. Based on this per-frame idea, 3DGStream \cite{sun20243dgstream} pioneers dynamic scene reconstruction via on-the-fly training, while subsequent works improve compression efficiency through techniques like rate-aware training \cite{hu20254dgcrateaware4dgaussian}, vector-quantized residuals \cite{girish2024queen}, and static-dynamic decomposition \cite{wu2025swift4d}.


\begin{figure}[t]
  \centering
   \includegraphics[width=1\linewidth]{imgs/teaser.pdf}

   \caption{ \textbf{Left:} Illustration of differences between previous optimization-based methods and our D-FCGS. \textbf{Right:} R-D curve and storage size comparison on N3V dataset.
    }
   \label{fig:teaser}
\end{figure}

% Motivation
Despite these advances, existing methods remain constrained by their coupled optimization of reconstruction and compression. By representing frame-wise Gaussian motions through specialized formats (e.g., neural networks), these approaches require per-scene optimization and meticulous hyperparameter tuning. This optimization-dependent paradigm creates two critical barriers to practical deployment: (1) it severely limits generalization to unseen scenes; (2) it hinders the development of standardized compression schemes that could enable widespread adoption of FVV. 

To address these challenges, we propose \textbf{D-FCGS}, a novel  feedforward compression framework for dynamic Gaussian Splatting. Our key insight is that temporal coherence in Gaussian point clouds can be efficiently modeled via the Group-of-Frames (GoF) representation, where inter-frame motions can be compressed in a feedforward and scene-agnostic manner. Specifically, we adopt the standard GS format as input and leverage sparse control points to efficiently extract motion tensors, optimal for both compression efficiency and computational performance. These motion tensors are then processed through our feedforward motion compression pipeline, which incorporates a tailored dual prior-aware entropy model to enhance probability estimation accuracy. Following decompression, the sparse motions are propagated across the entire Gaussian frame under the guidance of the control points. Finally, lightweight color refinement is applied to improve view-consistent fidelity.

Our D-FCGS is trained end-to-end on frame-wise GS sequences constructed from both real-world and synthetic multi-view video datasets. Once trained, it serves as \textbf{a general-purpose inter-frame compression codec}, requiring no scene-specific optimization or access to multi-view images during zero-shot inference. Extensive experiments show that D-FCGS exhibits strong generalization across diverse dynamic scenes and achieves state-of-the-art rate-distortion performance, surpassing 17$\times$ compression over 3DGStream while maintaining comparable fidelity.

Our contributions can be summarized as follows: 
\begin{itemize}
\item We present D-FCGS, a novel feedforward compression framework for dynamic 3DGS that enables zero-shot inter-frame coding of Gaussian sequences.
\item We propose a sparse motion representation with control points for I-P coding, coupled with a dual prior-aware entropy model for efficient compression, and develop a decoder with motion compensation and refinement for enhanced fidelity.
\item Experiments across various scenes show the effectiveness and robustness of D-FCGS, achieving over 17$\times$ compression over 3DGStream while preserving fidelity, outperforming most optimization-based methods.
\end{itemize}

\section{Related Work}

\subsection{3D Gaussian Splatting Compression} 

High storage demands of 3DGS have motivated compression efforts. Optimization-based compression methods include value-based and structure-based ones. The former involve the use of pruning \cite{girish2024eagles, ali2024trimming}, masking \cite{lee2024compact, wang2024end}, vector quantization \cite{niedermayr2024compressed, navaneet2023compact3d} and distillation \cite{fan2024lightgaussian} to reduce insignificant Gaussians or redundancy of Gaussian parameters. The latter utilize structural modeling such as anchors \cite{lu2024scaffold}, tri-planes \cite{lee2025compression} and 2D grids \cite{morgenstern2024compact} to address the sparsity and unorganized nature of Gaussian Splatting. When combined with learned entropy models \cite{balle2016end, cheng2020learned}, they achieve excellent rate-distortion performance \cite{chen2024hac, chen2025hac++, zhan2025cat}.

However, these optimization-based methods require per-scene fine-tuning, limiting their generalization. Recent works have introduced feedforward pipelines \cite{chen2024fast,  huang2025hierarchical, yang2025hybridgs} that decouples reconstruction and compression. These pre-trained codecs can directly compress arbitrary Gaussian clouds without multi-view supervision, offering greater practicality. We aim to extend this paradigm to dynamic GS compression.

\begin{figure*}[ht]
  \centering
   \includegraphics[width=1\linewidth]{imgs/overview.pdf}

   \caption{ \textbf{Overview of D-FCGS framework.} Our feedforward pipeline processes sequential Gaussian frames in GoF through three stages: (1) sparse motion extraction (\cref{sec:motion_extraction}), (2) feedforward motion compression (\cref{sec:motion_compression}), and (3) motion compensation and refinement (\cref{sec:frame_recon}). Once trained with rate-distortion loss, D-FCGS can infer on brand-new GS sequences in a zero-shot manner.}

   \label{fig:overview}
\end{figure*} 

\subsection{Dynamic GS and its Compression}

Dynamic GS methods can be categorized by their motion representation approaches. Implicit/explicit motion fields \cite{yang2024deformable, wu20244d, li2024spacetime, lin2024gaussian} and 4D Gaussian formulations \cite{duan20244d, lee2024fully} often face challenges in real-time streaming, variable resolutions, or long durations. In contrast, per-frame approaches \cite{luiten2024dynamic, sun20243dgstream, wang2024v, yan2025instant} demonstrate superior practicality through incremental frame updates, making them ideal for free-viewpoint video streaming. Our work focuses on this line of \textit{Per-Frame Gaussian} methods.

For compression of \textit{Per-Frame Gaussian}, HiCoM \cite{gao2024hicom} employs hierarchical grid-wise motion representation and QUEEN \cite{girish2024queen} leverages a latent-decoder for quantization of attribute residuals. 4DGC \cite{hu20254dgcrateaware4dgaussian} instills rate-aware training into frame reconstruction, achieving decent compression. While effective, they fundamentally couple reconstruction with compression, requiring scene-specific tuning. Our work breaks this constraint by introducing a novel optimization-free compression framework for dynamic GS, enabling zero-shot compression to arbitrary 4D scenes.


\section{Preliminary}
Our feedforward compression method takes frame-wise standard 3DGS as input, denoted as Gaussian frames, and each frame is pre-optimized from multi-view images at the respective timestamp.  Within a given Gaussian frame, each Gaussian primitive is characterized by: (1) \textbf{geometry parameters} including position $\boldsymbol{\mu} \in \mathbb{R}^3$ and covariance matrix $\boldsymbol{\Sigma} \in \mathbb{R}^{3\times3}$; (2) \textbf{appearance parameters} including opacity $o \in \mathbb{R}^1$ and SH-based color $c_{SH} \in \mathbb{R}^{48}$. The covariance matrix can be further represented as $\boldsymbol{\Sigma } = \boldsymbol{R}\boldsymbol{S}\boldsymbol{S}^T\boldsymbol{R}^T$, where $\boldsymbol{R} \in \mathbb{R}^{3\times3}$ is the rotation matrix parameterized by the quaternion $\boldsymbol{q} \in \mathbb{R}^4$,  and the scale matrix $\boldsymbol{S} \in \mathbb{R}^{3\times3}$ is a diagonal matrix with elements $\boldsymbol{s} \in \mathbb{R}^3$. The geometry of a Gaussian primitive can be formulated as:
\begin{equation}
  G(\boldsymbol{x}) = e^{-\frac{1}{2}(\boldsymbol{x} - \boldsymbol{\mu})^T \boldsymbol{\Sigma}^{-1} (\boldsymbol{x} - \boldsymbol{\mu})}, 
\end{equation}
where $\boldsymbol{x} \in \mathbb{R}^3$ is any random 3D location within the scene. 

Given a viewpoint, 3D Gaussians are projected into a 2D plane, and the color of a pixel $C \in \mathbb{R}^3$ is derived by alpha-blending of overlapping 2D Gaussians:
\begin{equation}
  C = \sum_i c_i \alpha_i \prod_{j=1}^{i-1} (1-\alpha_j), 
\end{equation}
where $c_i \in \mathbb{R}^3$ is the view-dependent color calculated from SH-based color $c_{SH}$, and $\alpha_i \in \mathbb{R}^1$ is the blending weight derived from opacity $o$.

With the differentiable rasterizer, training can be supervised by 2D images from varied views in an end-to-end way. The rendering loss of vanilla 3DGS is:
\begin{equation}
  L_{\text{render}} = \lambda L_{\text{D-SSIM}} + (1-\lambda) L_1.
\end{equation}


\section{Method}
\subsection{Overview}

As shown in \cref{fig:overview}, our D-FCGS framework comprises three key components: (1) sparse motion extraction (\cref{sec:motion_extraction}), (2) feedforward motion compression (\cref{sec:motion_compression}) and (3) motion compensation and refinement (\cref{sec:frame_recon}). Training and inference procedures will be mentioned in \cref{sec:training}.

Given two adjacent Gaussian frames $\boldsymbol{x_t}$ and $\boldsymbol{\hat{x}_{t-1}}$, we first sample sparse control points from dense Gaussians and extract motion features. These motion tensors are compressed using our dual prior-aware entropy model for efficient rate estimation. During decoding, we reconstruct $\boldsymbol{\hat{x}_t}$ via control-point-guided motion compensation and color refinement, while storing  it into the buffer for future reference.

\subsection{Sparse Motion Extraction via Control Points}
\label{sec:motion_extraction}
Building on observations that most Gaussians exhibit local motion coherence, we propose an efficient sparse representation. Inspired by previous optimization-based 4D reconstruction methods \cite{he2024s4d, yan2025instant, huang2024sc}, we employ Farthest Point Sampling (FPS) to select $N^c = \frac{N}{M}$ control points from the $N$ Gaussian primitives, and derive corresponding geometry parameters, formulated as:
\begin{equation}
    \boldsymbol{\mu^c} = FPS(\{\mu_i\}_{i\in N}, \frac{N}{M}),
\end{equation}
\begin{equation}
    \boldsymbol{x^{c}} = Index(\boldsymbol{x}, \boldsymbol{\mu^c}),
\end{equation}
where $M$ denotes the downscale factor, and $\boldsymbol{x^{c}}$  ($\boldsymbol{x^c} = \{\boldsymbol{\mu^c}, \boldsymbol{q^c}\}$) represent geometry parameters of control points. This process efficiently reduces storage demand and processing costs, while preserving motion characteristics.


After control point sampling, we derive the native sparse geometry attributes of the current frame and reference Gaussian frame, denoted as $\boldsymbol{x^c_t}$ and $\boldsymbol{\hat{x}^c_{t-1}}$. We encode these attributes using frequency encoding \cite{mildenhall2021nerf} and MLP projection:
\begin{equation}
    \boldsymbol{y^c_{t}} = MLP(FreqEnc\{\boldsymbol{x^c_t}\}),
\end{equation}
\begin{equation}
    \boldsymbol{\hat{y}^c_{t-1}} = MLP(FreqEnc\{\boldsymbol{\hat{x}^c_{t-1}}\}).
\end{equation}

Motion tensors are then derived in the feature domain: 
\begin{equation}
    \boldsymbol{m^c_{t}}   = Converter (\boldsymbol{y^c_{t}} - \boldsymbol{\hat{y}^c_{t-1})},
\end{equation}
where $\boldsymbol{m^c_{t}}$ denotes the motion tensors at time $t$, which keeps the same dimension as $\boldsymbol{x^c_t}$ and $\boldsymbol{\hat{x}^c_{t-1}}$. $Converter$ is realized by MLP at the feature level.

\subsection{Feedforward Motion Compression}
\label{sec:motion_compression}


\paragraph{\textbf{End-to-End Motion Compression.}}
The obtained sparse motion tensors $\boldsymbol{m^c_{t}}$ are then fed into our end-to-end compression module. This process begins with data encoding, followed by differentiable quantization simulated by additive uniform noise \cite{balle2016end} to enable gradient backpropagation :
\begin{equation}
    \begin{aligned}
    \boldsymbol{\hat{y}^m_t} =Q(\boldsymbol{y^m_t}) 
    & = \boldsymbol{y^m_t} + \mathcal{U}(-\frac{q'}{2}, \frac{q'}{2}),  \text{ for training} \\
    & = Round(\frac{\boldsymbol{y^m_t}}{q'}) \cdot q', \text{ for testing} \\
    \end{aligned}
\end{equation}
where $q'\in \mathbb{R}^1$ is the quantization step size, and $\boldsymbol{y^m_t}$, $\boldsymbol{\hat{y}^m_t}$ denote the encoded latent motion before and after quantization, respectively. Next to that, arithmetic coding (AC) converts the quantized data into a compact bitstream for efficient transmission and storage. On the decoder side, the bitstream is decompressed back into motion tensors using arithmetic decoding (AD) for frame reconstruction.

\paragraph{\textbf{Dual Prior-Aware Entropy Model.}}
According to Shannon’s theory \cite{shannon1948mathematical}, the cross-entropy between the estimated and true latent distributions provides a tight lower bound on the achievable bitrate:
\begin{equation}
    R(\boldsymbol{\hat{y}^m_t}) \ge \mathbb{E}_{\boldsymbol{\hat{y}^m_t} \sim q_{\boldsymbol{\hat{y}^m_t}}}[-\log _2p_{\boldsymbol{\hat{y}^m_t}}(\boldsymbol{\hat{y}^m_t})],
\end{equation}
where $p_{\boldsymbol{\hat{y}^m_t}}$ and $q_{\boldsymbol{\hat{y}^m_t}}$ are respectively estimated and true probability mass functions (PMFs) of the quantized latent codes $\boldsymbol{\hat{y}^m_t}$. Since arithmetic coding can achieve a bitrate close to this bound, our goal is to devise an entropy model that accurately estimates $p_{\boldsymbol{\hat{y}^m_t}}$.

\begin{figure}
    \centering
    \includegraphics[width=1\linewidth]{imgs/entropy.pdf}
    \caption{\textbf{Illustration of the dual prior-aware entropy model.} Spatial-temporal context priors extracted via multi-scale hashgrids and hyperpriors generated through the factorized model are fused by a lightweight fusion network for bitrate estimation.}

    \label{fig:entropy}
\end{figure}

\cref{fig:entropy} shows the proposed dual prior-aware entropy model that combines hyperprior and spatial-temporal context prior for precise distribution estimation. 
Following \cite{balle2018variational}, we use a factorized model to learn the hyperprior and estimate its PMF $p(\boldsymbol{\hat{z}^m_t})$, which is common in deep image and video compression. However, for GS-based FVV, the latent codes also exhibit strong 3D-spatial and temporal correlations. 
Assuming adjacent Gaussian frames share similar features, we strive to extract spatial-temporal priors from the reference frame $\boldsymbol{\hat{x}_{t-1}}$ through a multi-resolution hash grid encoding scheme. Specifically, at each grid intersection, we store learnable feature vectors that capture position-specific information. For each Gaussian positioned at $\boldsymbol{\mu_t}$, we retrieve multi-scale features by performing tri-linear interpolation across different levels $G^l$ of the voxel grid, where $l$ indicates the resolution level. The interpolated features from all levels are then concatenated and processed by a lightweight MLP to generate comprehensive positional contexts:
\begin{equation}
    context = MLP( \bigcup_{l=1}^L Interp(\boldsymbol{\mu_t}, G^l)),
\end{equation}
where $Interp(\cdot)$ denotes the grid interpolation operation. These positional contexts are subsequently combined with the transformed appearance parameters as the final spatial-temporal priors. The prior fusion network is then to integrate these spatial-temporal priors with the hyperpriors, estimating the mean $\boldsymbol{\mu^m_t}$ and scale $\boldsymbol{\sigma^m_t}$ of the latent code distribution (assumed as normal distribution). The probability mass for each quantized latent value $\hat{y}^m_{t,i}$ is computed by integrating the PMF over the quantization bin:
\begin{equation}
    p(\hat{y}^m_{t, i}) = \int_{\hat{y}^m_{t, i} - \frac{q'}{2}} ^ {\hat{y}^m_{t, i} + \frac{q'}{2}} \mathcal{N}(y|\mu^m_{t, i}, \sigma^m_{t, i}) dy,
\end{equation}
where $i$ corresponds to the index of a certain control point. The rate loss combines contributions from both the motion and hyper latent distributions:
\begin{equation}
    L_{\text{rate}} = \frac{1}{N^c} \sum_{i=1}^{N^c} (-\log_2( p(\hat{y}^m_{t, i}))  -\log_2(p(\hat{z}^m_{t, i}))),
\end{equation}
where $N^c = \frac{N}{M}$ is the total number of control points.

 
\subsection{Motion Compensation and Refinement}
\label{sec:frame_recon}


\paragraph{\textbf{Control Point Guided Motion Compensation.}}
The decoded control point motions $\boldsymbol{\hat{m}^c_t}$ are propagated to the entire Gaussian set in a distance-aware compensation way. For the $i^{th}$ control point, we first identify its $K$-nearest Gaussians $\mathcal{K}(i)$ via KNN search. The motion vectors are then distributed to neighboring Gaussians using an exponentially decaying weight function based on spatial distance. The closer the distance is, the more influence we consider the control point can exert on the motion of this neighboring point. The motion that $i^{th}$ control point assigns to its $j^{th}$ neighbor can be written as:
\begin{equation}
m_{i,j} = \frac{e^{-d_{i, j}} m_i}{\sum_{k\in \mathcal{K}(i)}e^{-d_{i, k}}},
\end{equation}
where $d_{i, j}$ represents the Euclidean distance between the $i^{th}$ control point and its $j^{th}$ neighbor. The aggregated motion vectors are finally applied to adjust the geometry parameters of the reference frame via addition. 

This approach provides two key benefits: (1) enhanced motion estimation accuracy through localized correlation modeling that suppresses error propagation, and (2) inherent parallelizability due to the independent processing of control points, ensuring computational efficiency. More discussions on frame matching during motion compensation can be seen in the supplementary.


\paragraph{\textbf{Color Refinement.}}
For image and video compression, post-compression refinement plays a critical role in mitigating visual artifacts such as color banding and blurring. 
%
In our GS-based framework, we specifically target refinement at SH coefficients $c_{SH}$ while preserving sensitive geometry and opacity parameters \cite{chen2024fast, girish2024queen, papantonakis2024reducing}.
%
The spatial-temporal priors from the entropy model module are repurposed to predict color residuals $\Delta c_{SH}$, which are dynamically added to $c_{SH}$ during decoding without additional storage. This \textbf{on-the-fly} refinement is fully differentiable, allowing gradient backpropagation through the entropy model for joint rate-distortion optimization.


\subsection{Training and Inference Pipeline of D-FCGS}
\label{sec:training}

\paragraph{\textbf{Training Process and Loss.}}
Following established practices \cite{wang2023neural, zheng2024hpc, zheng2024jointrf}, we adopt a Group-of-Frames (GoF) paradigm with the structure:
\begin{equation}
\scalebox{0.8}{$\scriptstyle
\underbrace{(I-P\!-\!\cdots\!-P)}_{\text{GoF}_1(L)} \quad
\underbrace{(I-P\!-\!\cdots\!-P)}_{\text{GoF}_2(L)} \quad
\cdots \quad
\underbrace{(I-P\!-\!\cdots\!-P)}_{\text{GoF}_k(L)},$}
\end{equation}
where $\text{GoF}_k(L)$ denotes the $k^{th}$ group containing one intra-coded (I) frame followed by $L-1$ predictively-coded (P) frames. The model is trained end-to-end with a composite loss function:
\begin{equation}
L_{\text{total}} = L_{\text{render}} + \lambda_{\text{size}} L_{\text{rate}},
\end{equation}
where $L_{\text{render}}$ replicates the original 3D Gaussian Splatting objective, and $\lambda_{\text{size}}$ controls the rate-distortion trade-off.

\paragraph{\textbf{Encoding and Decoding Process.}}
The encoder extracts sparse motion from control points and compresses motion tensors $\boldsymbol{\hat{y}^m_t}$ and hyperpriors $\boldsymbol{\hat{z}^m_t}$ via arithmetic coding. During decoding, the system first reconstructs $\boldsymbol{\hat{z}^m_t}$, then combines it with spatial-temporal priors to estimate distribution parameters ($\boldsymbol{\sigma^m_t}$, $\boldsymbol{\mu^m_t}$) for decoding $\boldsymbol{\hat{y}^m_t}$. These decoded motion features enable subsequent motion compensation.


\section{Experiments}
We rigorously evaluate our D-FCGS model by addressing three critical research questions:
\begin{itemize}
\item \textbf{Generalization and Effectiveness:} Can D-FCGS generalize across new scenes from widely used datasets while achieving competitive rate-distortion performance compared to optimization-based methods? (\cref{sec:Q1})
\item \textbf{Robustness and Stability:} How does D-FCGS perform under diverse and high-dynamic scenes, and is the system stable under varying hyperparameters? (\cref{sec:Q2})
\item \textbf{Module Efficacy:} Do individual modules (e.g. control points) contribute meaningfully? (\cref{sec:Q3})
\end{itemize}

\subsection{Experimental Setup}
\label{sec:exp}

\paragraph{\textbf{Datasets and Implementation Details.}} We derive sequential Gaussian frames from six multi-view video datasets: (1) \textit{N3V} \cite{li2022neural} (2) \textit{MeetRoom} \cite{li2022streaming} (3) \textit{WideRange4D} \cite{yang2025widerange4d} (4) \textit{Google Immersive} \cite{broxton2020immersive} (5) \textit{Self-Cap} \cite{xu2024longvolcap} (6) \textit{VRU} \cite{wu2025swift4d}. For training, we use 3 scenes from \textit{MeetRoom} and 28 scenes from \textit{WideRange4D}. For evaluation, we reserve the "discussion" scene from \textit{MeetRoom} for in-domain testing and six scenes from \textit{N3V} for out-of-domain benchmarking. Additional robustness tests are performed on diverse and high-dynamic scenes (see \cref{sec:Q2}). All experiments are conducted on NVIDIA RTX 4090 GPU. Key hyperparameters include: (1) the downscale factor $M$ = 70, (2) $K$ = 30 for KNN, (3) quantization step $q'$ = 1, (4) GoF size $L=10$ during inference, and (4) $\lambda_{\text{size}} $ = 1e-3. Further details of datasets and implementation are provided in the supplementary material.

\begin{figure*}[t]
  \centering
   \includegraphics[width=1\linewidth]{imgs/subjective.pdf}
\caption{\textbf{Qualitative comparison of rendering results.} Visual comparisons between (1) Ground Truth, (2) 3DGStream, and (3) our D-FCGS. D-FCGS significantly reduces storage size while maintaining comparable high fidelity to 3DGStream.}
   \label{fig:showcase}
\end{figure*}




\paragraph{\textbf{Evaluation Metrics.}} We evaluate D-FCGS using three kinds of metrics: (1) PSNR and SSIM \cite{wang2004image} for reconstruction quality, (2) compressed size (MB/frame) for rate efficiency, and (3) rendering speed (FPS) plus encoding/decoding time (sec) for computational performance.


\begin{table}[tb]
    \centering
    \small
    \setlength{\tabcolsep}{2pt}
        \begin{tabular}{c|cccc|c}
            \toprule
            Method & \begin{tabular}[c]{@{}c@{}}PSNR\\ (dB) ↑\end{tabular} & SSIM ↑ &  \begin{tabular}[c]{@{}c@{}}Size\\(MB) ↓\end{tabular} & \begin{tabular}[c]{@{}c@{}}Render\\ (FPS) ↑\end{tabular}   & \begin{tabular}[c]{@{}c@{}}Feedforward\\ Compression\end{tabular} \\ 
            \midrule
            K-Planes & 31.63& 0.920 & 1.0 & 0.15 &\ding{55} \\ 
            HyperReel & 31.10 & 0.931 & 1.7 & 16.7 &\ding{55} \\ 
            NeRFPlayer & 30.69 & 0.931 & 18.4& 0.05 & \ding{55} \\ 
            StreamRF & 30.61 & 0.930 & 7.6 & 8.3 & \ding{55} \\ 
            ReRF & 29.71 & 0.918 & 0.77 & 2.0 &\ding{55} \\ 
            TeTriRF & 30.65 & 0.931 & 0.76 & 2.7 & \ding{55} \\ 
            D-3DG & 30.67 & 0.931 & 9.2 & \textbf{460} & \ding{55} \\ 
            3DGStream* & \textbf{32.20} & \textbf{0.953} & 7.75 & 215 & \ding{55} \\ 
            HiCoM & 31.17 & - & 0.70 & \underline{274} & \ding{55} \\ 
            QUEEN & \underline{32.19} & 0.946 & 0.75 & 248 & \ding{55} \\ 
            4DGC & 31.58 & 0.943 & \underline{0.49} & 168 & \ding{55} \\              
            \textbf{D-FCGS (ours)} &  31.91 & \underline{0.952} & \textbf{0.46} & 215 &  \ding{51}\\ \bottomrule
        \end{tabular}
    \label{tab:n3v_result}
    \caption{
    Quantitative results on N3V \cite{li2022neural} dataset, averaged over 300 frames across six scenes. * denotes results reproduced by our implementation. \textbf{Bold} and \underline{underlined} values indicate the best and second-best performance, respectively. Detailed per-scene results are reported in supplementary material.
    }
\end{table}


\begin{table}[t]
    \centering
    \small
    \setlength{\tabcolsep}{7pt}
    \begin{tabular}{c|cccc}
        \toprule
        Method & \begin{tabular}[c]{@{}c@{}}PSNR\\ (dB) ↑\end{tabular} & SSIM ↑ &  \begin{tabular}[c]{@{}c@{}}Size\\(MB) ↓\end{tabular} & \begin{tabular}[c]{@{}c@{}}Render\\ (FPS) ↑\end{tabular}  \\ 
        \midrule
        StreamRF & 26.71 & 0.913 & 8.23 & 10  \\ 
        ReRF & 26.43 & 0.911 & 0.63 & 2.9 \\ 
        TeTriRF & 27.37 & 0.917 & 0.61 & 3.8  \\ 
        3DGStream* & \textbf{31.74} & \textbf{0.957} & 7.66 & \textbf{288} \\ 
        HiCoM & 29.61 & - & \underline{0.40}& \underline{284} \\ 
        4DGC & 28.08 & 0.922 & 0.42 & 213 \\ 
        \textbf{D-FCGS (ours)} & \underline{30.97} & \underline{0.950} & \textbf{0.38} & \textbf{288} \\ 
        \bottomrule
    \end{tabular}
    \label{tab:meetroom_result}
    \caption{
    Quantitative results on MeetRoom \cite{li2022streaming} dataset, averaged over 300 frames.
    }
\end{table}



\begin{table*}[t]
  \small
  \centering
  \label{tab:robustness}
  \begin{tabular}{c|ccc|ccc|ccc}
    \toprule
      & \multicolumn{3}{c|}{\textbf{Google Immersive}}
 & \multicolumn{3}{c|}{\textbf{Self-Cap}}  & \multicolumn{3}{c}{\textbf{VRU}}  \\
Method & PSNR(dB)↑ & SSIM↑ & Size(MB)↓ & PSNR(dB)↑ & SSIM↑ & Size(MB)↓ & PSNR(dB)↑ & SSIM↑ & Size(MB)↓  \\
    \midrule
     3DGStream* &25.84& 0.872& 7.6 & 26.43 & 0.858 & 7.6 & 26.18 & 0.892 & 7.6 \\
     D-FCGS(ours) & 25.58& 0.879 &0.34& 26.00 & 0.854 & 0.90 & 25.13 & 0.876 & 0.43\\
    \bottomrule
\end{tabular}
  \caption{Robustness test on diverse scenes from extensive datasets. Averaged PSNR, SSIM and P-frame size are reported. Per-scene results are provided in the appendix.}
\end{table*}

\begin{table}[h]
  \centering
  \label{tab:hyperparameters_robustness}
   \small
   \setlength{\tabcolsep}{2pt}
   
  \begin{tabular}{c|ccccccc}
    \toprule
    $K$ & 30 & 60 & 120 & 30 & 30 & 30 & 30 \\

    $M$   & 70 & 70 & 70 & 140 & 280 & 70 & 70 \\

    $L$   & 300 & 300 & 300 & 300 & 300 & 32 & 4\\
    \midrule
    Size↓ & 8.9 & 9.0 & 9.0 & 4.0 & 2.0 & 9.0 & 9.2\\
     PSNR↑  & 33.596 & 33.596 & 33.596 & 33.599 & 33.599 & 33.604 & 33.598\\
    \bottomrule
  \end{tabular}
    \caption{Robustness test on hyperparameter settings ($K$ for KNN, downscale factor $M$, and GoF size $L$). Averaged P-frame size ($\times$0.1 MB) and PSNR (dB) are reported on "cut\_beef" scene.}
\end{table}

\begin{table}[t]
  \centering
  \label{tab:time}
  \small
  \setlength{\tabcolsep}{4pt}
  \begin{tabular}{c|ccc}
    \toprule
    Method&Encoding(sec) & Decoding(sec)  & Total(sec) \\
    \midrule
     proposed &0.61& 0.72& 1.33 \\
     w/o control points &1.33& 2.88& 4.21 \\
    \bottomrule
    \end{tabular}
  \caption{Average encoding and decoding time for P-frames.}
\end{table}

\subsection{Benchmark Comparison}
\label{sec:Q1}

\paragraph{\textbf{Benchmark Methods.}} As the first feedforward inter-frame codec for Gaussian point clouds (to our knowledge), D-FCGS faces unique evaluation challenges due to the absence of directly comparable optimization-free approaches. Our analysis therefore encompasses two categories of benchmark methods in optimization-based dynamic scene reconstruction: (1) NeRF-based techniques including K-Planes \cite{fridovich2023k}, HyperReel \cite{attal2023hyperreel}, NeRFPlayer \cite{song2023nerfplayer}, StreamRF \cite{li2022streaming}, ReRF \cite{wang2023neural}, and TeTriRF \cite{wu2024tetrirf}; and (2) \textit{Per-Frame Gaussian} methods including D-3DG \cite{luiten2024dynamic}, 3DGStream \cite{sun20243dgstream}, HiCoM \cite{gao2024hicom}, QUEEN \cite{girish2024queen}, and 4DGC \cite{hu20254dgcrateaware4dgaussian}. While this comparison inherently favors optimization-based approaches  that benefit from scene-specific tuning, it provides critical insights into D-FCGS's performance relative to current paradigms in dynamic scene compression. 

\begin{figure*}[t]
  \centering
   \includegraphics[width=1\linewidth]{imgs/rdcurve.pdf}
   % \includegraphics[width=1\linewidth]{imgs/ablation-rdcurve.pdf}
   \caption{Rate-Distortion comparison of the proposed method and ablations.}
   \label{fig:aba_rd}
   % \vspace{-1em}
\end{figure*}

% \begin{figure}[t]
%     \centering
%     \includegraphics[width=0.95\linewidth]{imgs/perframe.pdf}
%     \caption{
%     Quality fluctuations on "flame steak" scene.
% }
%     \label{fig:trend}
% \end{figure}

\paragraph{\textbf{Quantitative Results.}} 


 Our compression results demonstrate significant improvements over existing methods. As shown in \cref{tab:n3v_result} and \cref{tab:meetroom_result}, D-FCGS achieves remarkable average sizes of \textbf{0.46MB} (\textit{N3V}) and \textbf{0.38MB} (\textit{MeetRoom}) per frame, namely a \textbf{17×} reduction compared to 3DGStream's 7.75MB and 7.66MB. Note that our codec is designed for inter-frame compression, and our P-frame compression ratio exceeds 32× in most cases (\cref{fig:teaser}). While I-frames could be further compressed using existing static methods (e.g., FCGS \cite{chen2024fast}), we maintain their original size for fair comparison. Additionally, the entire encoding/decoding pipeline operates efficiently, with both processes completing in under 1 second (\cref{tab:time}).

\paragraph{\textbf{Qualitative Results.}} We visualize the qualitative comparisons with 3DGStream \cite{sun20243dgstream} in \cref{fig:showcase}, presenting results on "flame steak", "flame salmon" (\textit{N3V} dataset) and "discussion" (\textit{MeetRoom} dataset). From scene details, we can tell that D-FCGS achieves comparable fidelity to 3DGStream, even better in some cases.


\subsection{Robustness Test}
\label{sec:Q2}
\paragraph{Robustness to Diverse Scenes.}
We evaluate D-FCGS on three additional datasets: \textit{Google Immersive}, \textit{Self-Cap}, and \textit{VRU}, covering scenarios such as indoor painting, cave exploration, and basketball games. As summarized in \cref{tab:robustness}, for Gaussian sequences reconstructed by 3DGStream, our approach achieves unprecedented P-frame compression while maintaining near-identical SSIM and small PSNR drop ($<$ 0.5 dB for \textit{Google Immersive} and \textit{Self-Cap}). While PSNR gaps exist for \textit{VRU} (full of blurry motions), our method offers a more practical and efficient pipeline for real-world applications, sacrificing acceptable visual fidelity for great improvements in bitrates. \cref{fig:showcase} highlights visual results on the outdoor "cave" (\textit{Google Immersive}) and high-dynamic "bar" (\textit{Self-Cap}). 

% Per-scene results and additional visualizations are provided in the supplementary.


\paragraph{Robustness to Hyperparameter Settings.} Here, we apply varied parameter configurations (KNN neighborhood size $K$, downscale factor $M$, and GoF size $L$) to the "cut beef" scene (\textit{N3V}). As presented in \cref{tab:hyperparameters_robustness}, D-FCGS maintains stable rendering quality across most parameter combinations, showing strong parametric robustness. Notably, increasing the downscale factor $M$ effectively reduces control point counts, leading to a linear decrease in P-frame compression size, which aligns with expectations.


\subsection{Ablation Study}
\label{sec:Q3}

In this section, we evaluate key components of D-FCGS through systematic ablation studies.
% and present the quality fluctuation during inter-frame encoding/decoding process.



\paragraph{\textbf{Effect of Control Points.}} The sparse motion representation via control points forms the foundation of our efficient compression pipeline. Removing control points and predicting motions for all Gaussians increases storage substantially (\cref{fig:aba_rd}) and slows encoding/decoding by 3.2× (\cref{tab:time}). This validates our sparse motion representation's efficiency in rate saving and computational cost.


\paragraph{\textbf{Effect of Dual Prior-Aware Entropy Model.}} 
 Our entropy model employs a novel dual-prior architecture. While hyperpriors effectively capture global dependencies  \cite{balle2018variational}, our key innovation lies in the hash-grid-based spatial-temporal prior that models local correlations. Removing the spatial-temporal prior branch leads to noticeable R-D performance degradation (\cref{fig:aba_rd}), confirming its importance to optimal compression.



\paragraph{\textbf{Effect of Color Refinement.}} 

Our online color refinement module improves rendering quality (PSNR) by  0.1 $\sim$ 0.5 dB (\cref{fig:aba_rd}), while requiring no additional storage and  negligible decoding overhead. Per-scene results on \textit{N3V} dataset are shown in the supplementary.


% \paragraph{\textbf{Analysis of Quality Fluctuations.}} As depicted in \cref{fig:trend}, D-FCGS maintains remarkably stable and marginal decompression quality gaps to the original 3DGStream-generated Gaussians  throughout the entire sequence, even when processing large GoFs ($L$ = 300). This demonstrates our method's effectiveness in preserving visual fidelity across extended temporal segments.

\section{Conclusion}


In this paper, we propose \textbf{F}eedforward \textbf{C}ompression of \textbf{D}ynamic \textbf{G}aussian \textbf{S}platting (D-FCGS), a novel feedforward framework for zero-shot dynamic Gaussian sequence compression. 
%
Our contributions are threefold. 
%
First, we adopt the I-P coding profile for standard GS compression and introduce sparse inter-frame motion extraction via control points. 
%
Second, we present an end-to-end motion compression framework with a dual prior-aware entropy model, fully leveraging hyperpriors and spatial-temporal context to improve rate estimation. 
%
Third, control-point-guided motion compensation is combined with a color refinement network to guarantee high-fidelity and view-consistent reconstruction. 
%
Experimental results show that \emph{D-FCGS} achieves superior compression efficiency (over 17× compression) on two benchmark datasets (MeetRoom and N3V) and remarkable robustness across diverse scenes from extensive datasets, significantly enhancing transmission and storage efficiency for free-viewpoint video applications.



\section*{Acknowledgments}
This work was partly supported by the NSFC62431015, Science and Technology Commission of Shanghai Municipality No.24511106200, the Fundamental Research Funds for the Central Universities, Shanghai Key Laboratory of Digital Media Processing and Transmission under Grant 22DZ2229005, 111 project BP0719010, Okawa Research Grant and Explore-X Research Fund.

\bibliography{aaai2026}


\end{document}
